% Continuous worked model for detailed chapters 03--04.
\section{综合例题：同一区间算子的定义域、谱、演化与计算误差}
\label{sec:ms-interval-spectral-comprehensive}

同一个微分表达式 $-u''$，配上不同边界条件以后，会给出不同的谱和不同的长期演化。
要把\chapref{chap:math-hilbert} 与\chapref{chap:math-spectrum} 的一般结论用于这个问题，必须先确定函数空间，再检查伴随和边界域。
我们固定区间 $(0,\pi)$，并始终使用同一个初始函数
\[
 f(x)=x(\pi-x).
\]
我们要解决的完整问题是：表达式 $-u''$ 怎样成为自伴算子，边界条件怎样改变它的谱，怎样求静态响应与含时解，以及只保留有限个模态时误差究竟有多大。
静态响应和演化都从选定的自伴算子出发，有限模态误差则由同一组谱系数计算。

\begin{exbox}[函数空间、初态范数与总量泛函]
取复 Hilbert 空间
\[
 \mathcal H=L^2(0,\pi),\qquad
 \langle u,v\rangle=\int_0^\pi u(x)\overline{v(x)}\,\dd x.
\]
这里沿用前两章对第一变量线性的内积约定。元素是几乎处处相等的函数的等价类；因此现在还不能给一般元素规定端点值。

\textbf{解。} 初态的范数可以直接积分：
\begin{align*}
 \|f\|^2
 &=\int_0^\pi (\pi^2x^2-2\pi x^3+x^4)\,\dd x\\
 &=\pi^5\left(\frac13-\frac12+\frac15\right)
 =\frac{\pi^5}{30}.
\end{align*}
若要把它作为归一化态，使用 $\widehat f=\sqrt{30}\,\pi^{-5/2}f$；下面保留 $f$，以便静态解保持简单。

空间的完备性在本题中保证模态级数的极限确实仍是一个 $L^2$ 元素。
具体地，对任意 Cauchy 序列 $u_j$，可取子列满足
\[
 \|u_{j_{m+1}}-u_{j_m}\|\le 2^{-m}.
\]
由 Minkowski 不等式，部分和 $\sum_{m=1}^M|u_{j_{m+1}}-u_{j_m}|$ 的 $L^2$ 范数不超过 $1$。
单调收敛给出一个几乎处处有限的 $L^2$ 上界，故差分级数几乎处处绝对收敛，并且尾项的 $L^2$ 范数不超过 $\sum_{m\ge M}2^{-m}$。
其和定义 $u\in L^2$，子列在范数中趋于 $u$；原序列的 Cauchy 性再给出整个序列收敛。

总量泛函 $\ell(u)=\int_0^\pi u(x)\,\dd x$ 满足
\[
 \ell(u)=\langle u,1\rangle,\qquad
 |\ell(u)|\le\sqrt\pi\|u\|,\qquad \|\ell\|=\sqrt\pi.
\]
最后一个等号由 $u=\pi^{-1/2}$ 达到；Riesz 表示向量就是常函数 $1$。
本题初态的总量为
\[
 \ell(f)=\int_0^\pi x(\pi-x)\,\dd x=\frac{\pi^3}{6}.
\]
点值泛函则不能在 $L^2$ 上照搬：令 $b_m(x)=\max(1-mx,0)$，有 $b_m(0)=1$ 而 $\|b_m\|^2=1/(3m)$。
端点条件必须等到选定 Sobolev 定义域以后才有意义。
\end{exbox}

\begin{exbox}[最小算子、最大算子与伴随域]
上一段已经确定初态所在的空间。现在要判断 $-f''$ 是否有意义，并求这个表达式的最小、最大定义域。从
\[
 T_0u=-u'',\qquad \Dom(T_0)=C_c^\infty(0,\pi)
\]
出发。紧支撑光滑函数在 $L^2$ 中稠密，因此 $T_0^*$ 的定义有唯一性。

\textbf{解。} 若 $v\in\Dom(T_0^*)$，存在 $h\in L^2$ 使
\[
 \int_0^\pi(-u'')\overline v\,\dd x
 =\int_0^\pi u\overline h\,\dd x
 \quad\text{对所有 }u\in C_c^\infty.
\]
这正是分布意义的 $-v''=h$。
一维 Sobolev 正则性给 $v\in H^2(0,\pi)$；反过来，任意 $H^2$ 函数对紧支撑 $u$ 都满足该等式。
所以
\[
 \Dom(T_{\max})=H^2(0,\pi),\qquad
 T_0^*=T_{\max},\qquad T_{\max}v=-v''.
\]
这里的正则性也可直接看出：令 $w(x)=-\int_0^x(x-y)h(y)\,\dd y$，则 $w\in H^2$ 且 $-w''=h$。
于是 $(v-w)''=0$，分布意义下的余项是仿射函数，也属于 $H^2$。

在有限区间上，图范数
\[
 \|u\|_{T_{\max}}^2=\|u\|^2+\|u''\|^2
\]
与 $H^2$ 范数等价。
因此 $T_0$ 的图闭包就是紧支撑光滑函数在 $H^2$ 中的闭包：
\[
 \Dom(T_{\min})=H_0^2(0,\pi)
 =\{u\in H^2:u(0)=u(\pi)=u'(0)=u'(\pi)=0\}.
\]
四个迹均连续，故闭包必须满足这四个条件；反向包含由零延拓后的光滑逼近与端点截断得到。
于是
\[
 T_{\min}=\overline{T_0},\qquad
 T_{\min}^*=T_{\max},\qquad T_{\max}^*=T_{\min}.
\]
图闭性还可直接检查：若 $u_j\to u$ 且 $-u_j''\to h$ 在 $L^2$ 中成立，分布极限给 $-u''=h$，上述正则性给 $u\in H^2$；对最小域，图范数收敛进一步保留四个零迹。

$H_0^2$ 同时固定函数值和一阶导数，是最小域；$H^2\cap H_0^1$ 只固定函数值，将成为 Dirichlet 自伴域。
给定初态满足 $f\in H^2$、$f(0)=f(\pi)=0$，但 $f'(0)=\pi$、$f'(\pi)=-\pi$，所以它在最大域而不在最小域中；若仍只使用最小算子，就不能对这个初态开始演化。
\end{exbox}

\begin{exbox}[由边界型和亏量空间选择自伴扩张]
上一段得到最小域过小、最大域尚未固定边界。现在求允许初态 $f$ 的自伴实现，并列出改变边界后的其他实现。
对 $u,v\in H^2(0,\pi)$ 分部积分，得到
\begin{align*}
 \langle T_{\max}u,v\rangle-\langle u,T_{\max}v\rangle
 &=[u\overline{v'}-u'\overline v]_0^\pi\\
 &=\langle\Gamma_1u,\Gamma_0v\rangle_{\mathbb C^2}
   -\langle\Gamma_0u,\Gamma_1v\rangle_{\mathbb C^2},
\end{align*}
其中
\[
 \Gamma_0u=\binom{u(0)}{u(\pi)},\qquad
 \Gamma_1u=\binom{u'(0)}{-u'(\pi)}.
\]
两端导数的符号已经由 Green 式固定。
任给四个边界数据，都有一个三次多项式实现它们，故 $(\Gamma_0,\Gamma_1)$ 映满 $\mathbb C^2\oplus\mathbb C^2$。

\textbf{解。} 亏量方程为
\[
 -u''=\ii u\quad\text{或}\quad -u''=-\ii u.
\]
令 $a_+=(1-\ii)/\sqrt2$、$a_-=(1+\ii)/\sqrt2$，则
\[
 \mathcal N_+=\operatorname{span}\{e^{a_+x},e^{-a_+x}\},\qquad
 \mathcal N_-=\operatorname{span}\{e^{a_-x},e^{-a_-x}\}.
\]
有限区间上四个指数函数均平方可积，所以亏指数是 $(2,2)$。
与半直线相比，这里两个端点都需要边界数据。

全部自伴域可用 $U\in U(2)$ 写成
\[
 \Dom(A_U)=\{u\in H^2:
 \Gamma_1u+\ii\Gamma_0u
 =U(\Gamma_1u-\ii\Gamma_0u)\},\qquad A_Uu=-u''.
\]
边界型的消失等价于相应边界子空间各向同性；这条幺正关系选择一个二维最大各向同性子空间，故伴随域与自身相同。
$U=I$ 给 Dirichlet 条件，$U=-I$ 给 Neumann 条件。
也可把两端耦合：周期条件 $u(0)=u(\pi)$、$u'(0)=u'(\pi)$ 同样使边界型消失并给自伴域。

对最常用的分离边界条件，直接得到
\begin{align*}
 \Dom(A_D)&=\{u\in H^2:u(0)=u(\pi)=0\},\\
 \Dom(A_N)&=\{u\in H^2:u'(0)=u'(\pi)=0\},\\
 \Dom(A_R)&=\{u\in H^2:u'(0)=h_0u(0),\
 u'(\pi)=-h_\pi u(\pi)\},\quad h_0,h_\pi\in\mathbb R.
\end{align*}
例如，要从对称性升级到 $A_D=A_D^*$，让 $u\in\Dom(A_D)$ 的 $u'(0),u'(\pi)$ 独立变化；伴随边界型为零迫使 $v(0)=v(\pi)=0$。
对 Neumann 域让端点值独立变化，便得到 $v'(0)=v'(\pi)=0$。
对 Robin 域代入条件后，边界型分别含
\[
 u(\pi)\bigl(\overline{v'(\pi)}+h_\pi\overline{v(\pi)}\bigr),
 \qquad
 -u(0)\bigl(\overline{v'(0)}-h_0\overline{v(0)}\bigr),
\]
两端值仍可独立选取，故 $v$ 满足相同条件。
因此每个候选域都恰好等于自己的伴随域，三个实现均为自伴算子。
\end{exbox}

\begin{exbox}[闭二次型、Friedrichs 扩张与最低能量]
第三步的 Dirichlet 实现允许给定初态 $f$。选定这个实现后，计算初态的能量形式，并检查形式域和算子域的区别。

\textbf{解。} 在最初的紧支撑域上
\[
 \langle T_0u,u\rangle=\int_0^\pi|u'|^2\,\dd x.
\]
以 $\|u\|_q^2=\|u\|^2+\|u'\|^2$ 取闭包，得到闭形式
\[
 \Dom(q_D)=H_0^1(0,\pi),\qquad
 q_D(u,v)=\int_0^\pi u'\overline{v'}\,\dd x.
\]
它的关联算子域须额外满足：存在 $h\in L^2$，使 $q_D(u,v)=\langle h,v\rangle$ 对所有 $v\in H_0^1$ 成立。
先取紧支撑 $v$ 得 $-u''=h$，再用一维正则性得 $u\in H^2$；因此关联算子正是 $A_D$。
这给出了最小算子的 Friedrichs 扩张。
\[
 \boxed{\Dom(q_D)=H_0^1\,,\qquad
 \Dom(A_D)=H^2\cap H_0^1\,.}
\]
二者要求的一阶、二阶正则性不同。

对 Neumann 和实 Robin 条件，相应闭形式为
\begin{align*}
 \Dom(q_N)&=H^1,&q_N(u,v)&=\int u'\overline{v'}\,\dd x,\\
 \Dom(q_R)&=H^1,&q_R(u,v)&=\int u'\overline{v'}\,\dd x
 +h_0u(0)\overline{v(0)}+h_\pi u(\pi)\overline{v(\pi)}.
\end{align*}
迹不等式
\[
 |u(0)|^2+|u(\pi)|^2\le\varepsilon\|u'\|^2+C_\varepsilon\|u\|^2
\]
保证任意有限实 Robin 参数都给闭的下半有界形式；当参数非负时，形式本身非负。
其关联算子的 Robin 条件来自允许任意 $H^1$ 测试函数的边界项，不需要在形式域中预先加上导数边界条件。

最低能量可以用正弦系求出。记 $e_n(x)=\sqrt{2/\pi}\sin(nx)$；其完备性将在下一框验证。
对紧支撑光滑函数分部积分后使用 Parseval，再以形式范数取极限，得到
\[
 q_D(u,u)=\sum_{n\ge1}n^2|\langle u,e_n\rangle|^2
 \ge\sum_{n\ge1}|\langle u,e_n\rangle|^2=\|u\|^2.
\]
等号恰在 $u$ 是 $\sin x$ 的倍数时成立。
所以 $A_D\ge I$，最低谱值为 $1$；Neumann 的常函数却满足 $q_N(1,1)=0$，其最低谱值是 $0$。
\end{exbox}

\begin{exbox}[Dirichlet 本征系、谱测度与初态展开]
第四步的能量形式确定了正的 Dirichlet 实现。现在求 $A_Du=\lambda u$，构造谱测度，并把给定初态 $f$ 完整展开。

\textbf{解。} 非零 Dirichlet 解要求
\[
 u(x)=C\sin(\sqrt\lambda x),\qquad
 \sqrt\lambda=n\in\{1,2,\ldots\}.
\]
非正实 $\lambda$ 没有非零解。归一化后
\[
 \lambda_n=n^2,\qquad e_n(x)=\sqrt{\frac2\pi}\sin(nx),\qquad
 \langle e_n,e_m\rangle=\delta_{nm}.
\]
完备性来自奇延拓：若 $u$ 与全部 $e_n$ 正交，把它奇延拓至 $(-\pi,\pi)$，则正弦 Fourier 系数全为零，余弦与常数系数由奇性为零。
Fourier 完备性给延拓在 $L^2$ 中为零，故 $u=0$。

完备性和 Parseval 等式说明，映射
\[
 \mathcal Fu=(\langle u,e_n\rangle)_{n\ge1}
\]
是从 $\mathcal H$ 到 $\ell^2$ 的幺正映射。
在 $A_D$ 的定义域内，两次分部积分使 $A_D$ 变成各系数乘以 $n^2$。
反过来，若 $\sum n^4|\langle u,e_n\rangle|^2<\infty$，正弦部分和在图范数中收敛；$A_D$ 的闭性保证其极限仍在定义域内。因此
\[
 \Dom(A_D)=\left\{u:\sum_{n\ge1}n^4|\langle u,e_n\rangle|^2<\infty\right\},
 \qquad \mathcal F A_D\mathcal F^{-1}(c_n)=(n^2c_n).
\]
对 Borel 集 $B\subset\mathbb R$，谱投影实际作用为
\[
 E_D(B)u=\sum_{n^2\in B}\langle u,e_n\rangle e_n.
\]
于是本题谱测度是投影原子 $E_D(\{n^2\})$ 的和；谱积分和 Borel 函数演算成为
\[
 \varphi(A_D)u=\sum_{n\ge1}\varphi(n^2)\langle u,e_n\rangle e_n,
 \quad
 \Dom(\varphi(A_D))=\left\{u:\sum|\varphi(n^2)|^2|\langle u,e_n\rangle|^2<\infty\right\}.
\]

为求初态在每个能量上的权重，令 $I_n=\int_0^\pi f(x)\sin(nx)\,\dd x$。
第一次分部积分的边界项因 $f(0)=f(\pi)=0$ 消失，第二次的边界项因 $\sin(n\pi)=\sin0=0$ 消失。因 $f''=-2$，得到
\begin{align*}
 I_n&=\frac1n\int_0^\pi f'(x)\cos(nx)\,\dd x\\
 &=-\frac1{n^2}\int_0^\pi f''(x)\sin(nx)\,\dd x
 =\frac{2(1-(-1)^n)}{n^3}.
\end{align*}
故
\[
 c_n=\langle f,e_n\rangle
 =\begin{cases}4\sqrt{2/\pi}\,n^{-3},&n\text{ 为奇数},\\0,&n\text{ 为偶数},\end{cases}
 \qquad
 f(x)=\frac8\pi\sum_{\substack{n\ge1\\n\text{ 奇}}}\frac{\sin(nx)}{n^3}.
\]
Parseval 现在给一个可核对的数值恒等式：
\[
 \frac{\pi^5}{30}=\sum_n|c_n|^2
 =\frac{32}\pi\sum_{n\text{ 奇}}\frac1{n^6},
 \quad \sum_{n\text{ 奇}}n^{-6}=\frac{\pi^6}{960}.
\]
初态自身的标量谱测度于是完全确定：
\[
 \mu_f(B)=\langle E_D(B)f,f\rangle
 =\frac{32}\pi\sum_{\substack{n\text{ 奇}\\n^2\in B}}n^{-6}.
\]
若使用第一步的归一化态，测到能量 $n^2$ 的概率为 $960/(\pi^6n^6)$（奇数 $n$），偶数概率为零；所有概率之和为 $1$。
同样，$q_D(f,f)=\int_0^\pi(\pi-2x)^2\,\dd x=\pi^3/3$。
初态属于 $\Dom(A_D)$ 且 $A_Df=2$，但不属于 $\Dom(A_D^2)$：常函数 $2$ 不满足 Dirichlet 端点条件，且 $\sum n^8|c_n|^2$ 发散。
其图范数也有具体值 $\|f\|_{A_D}^2=\pi^5/30+4\pi$。
空间上是多项式，并不保证满足每一阶算子域的边界要求。
\end{exbox}

\begin{exbox}[改变边界条件：Robin 根与 Neumann 零模]
第五步已经得到 $f$ 的 Dirichlet 模态。现在保持空间、表达式和初态不变，只改变第三步中的自伴域，求新的谱与初态系数。

\textbf{解。} 对 Robin 条件及 $\lambda=k^2>0$，先用左端条件固定
\[
 u(x)=C\left(\cos(kx)+\frac{h_0}{k}\sin(kx)\right).
\]
右端条件化成
\[
 (h_0h_\pi-k^2)\sin(k\pi)
 +k(h_0+h_\pi)\cos(k\pi)=0.
\]
这些正根平方才是正本征值；不能把乘以 $k$ 后自动出现的 $k=0$ 当成本征值。
零能量须另外解 $u=ax+b$，其相容条件为
\[
 h_0+h_\pi+\pi h_0h_\pi=0.
\]
当 $h_0,h_\pi\ge0$ 时，仅 $h_0=h_\pi=0$ 有零模。
一般实参数还可能有负本征值，须将 $k=\ii\kappa$ 代入并求双曲函数方程；不能沿用 Dirichlet 的严格正性。

Neumann 情形完整本征系为
\[
 \eta_0=\pi^{-1/2},\quad \lambda_0=0;
 \qquad \eta_n=\sqrt{2/\pi}\cos(nx),\quad\lambda_n=n^2\quad(n\ge1).
\]
初态在这组基中的系数是
\begin{align*}
 d_0&=\langle f,\eta_0\rangle=\frac{\pi^{5/2}}6,\\
 d_n&=\sqrt{2/\pi}\int_0^\pi f(x)\cos(nx)\,\dd x
 =-\sqrt{2\pi}\,\frac{1+(-1)^n}{n^2}\quad(n\ge1).
\end{align*}
这次非零的是偶数余弦模态，外加一个常数模态。
$f\in H^1$ 是合法的 Neumann 能量初态，但 $f'(0)=\pi$、$f'(\pi)=-\pi$，故 $f\notin\Dom(A_N)$。
相应地，$\sum n^2|d_n|^2$ 有限，而 $\sum n^4|d_n|^2$ 发散。
改变边界域，同时改变了同一初态的谱系数和允许的时间正则性。
\end{exbox}

\begin{exbox}[预解式的谱表示与 Green 核]
前两段的本征系使静态响应可以按模态求和。现在把源记为 $g$，先求任意源的完整预解式，再在下一段代入本题的恒定源 $1$。
对 $z\notin\{n^2:n\ge1\}$，谱表示给
\[
 R_D(z)=(A_D-z)^{-1},\qquad
 R_D(z)g=\sum_{n\ge1}\frac{\langle g,e_n\rangle}{n^2-z}e_n,
 \quad \|R_D(z)\|=\frac1{\operatorname{dist}(z,\sigma(A_D))}.
\]
模态表示适合逐频率求解；若源在坐标空间直接给定，积分核能把源到响应的关系写得更直接。

\textbf{解。} 写 $z=k^2$，选左右边界解 $s_L(x)=\sin(kx)$、$s_R(x)=\sin(k(\pi-x))$。
其 Wronskian $s_Ls_R'-s_L's_R=-k\sin(k\pi)$。
在源点两侧，核必须分别正比于 $s_L$ 和 $s_R$，以满足两端的零边界值。
核在 $x=y$ 连续；把 $(-\partial_x^2-z)G=\delta_y$ 在源点附近积分，则得导数跳跃
$\partial_xG(y^+,y)-\partial_xG(y^-,y)=-1$。
用这两个条件固定系数，得到
\[
 G_D(z;x,y)=\frac{\sin(kx_<)\sin(k(\pi-x_>))}{k\sin(k\pi)},
 \qquad x_<=\min(x,y),\quad x_>=\max(x,y).
\]
它不依赖平方根分支的符号；$k=0$ 是可去点。
边界值为零、跳跃为 $-1$，所以 $(-\partial_x^2-z)G_D=\delta_y$。
因此
\[
 R_D(z)g(x)=\int_0^\pi G_D(z;x,y)g(y)\,\dd y,
 \qquad
 G_D(0;x,y)=\frac{x_<\,(\pi-x_>)}\pi.
\]
与谱式比较，得到两种完整表示之间的联系：
\[
 G_D(z;x,y)=\frac2\pi\sum_{n\ge1}\frac{\sin(nx)\sin(ny)}{n^2-z}.
\]
每个固定非谱 $z$ 的级数绝对且一致收敛，因其尾项由 $C_z/n^2$ 控制。

Neumann 的左右解改成余弦，核和谱式相应变为
\begin{align*}
 G_N(z;x,y)&=-\frac{\cos(kx_<)\cos(k(\pi-x_>))}{k\sin(k\pi)}\\
 &=-\frac1{\pi z}+\frac2\pi\sum_{n\ge1}\frac{\cos(nx)\cos(ny)}{n^2-z}.
\end{align*}
第一项是零模的极点，不能当作可去奇点删除。
去掉常数子空间后，逆的核为
\[
 G_N^\perp(x,y)=\frac{x^2+y^2}{2\pi}-\max(x,y)+\frac\pi3.
\]
在 $x\ne y$ 处，$-\partial_x^2G_N^\perp=-1/\pi$；在 $x=y$ 处，一阶导数跳跃为 $-1$，所以
$-\partial_x^2G_N^\perp=\delta_y-1/\pi$。
两端导数为零，且直接积分给 $\int_0^\pi G_N^\perp(x,y)\,\dd x=0$。
它只反演平均值为零的源；完整 Neumann 算子没有处处定义的逆。
\end{exbox}

\begin{exbox}[静态响应、奇异值分解与最佳有限秩逼近]
在第七步求出的预解式中代入 $z=0$、$g=1$，求 $A_Du=1$。同一问题既用微分方程、Green 核，也用谱计算，并检查所得响应和初态 $f$ 的关系。

\textbf{解。} 微分方程及边界条件给
\[
 u''=-1,\qquad u(0)=u(\pi)=0,
 \qquad u(x)=\frac{x(\pi-x)}2=\frac{f(x)}2.
\]
直接积分第七步的核也给
\begin{align*}
 \int_0^\pi G_D(0;x,y)\,\dd y
 &=\frac{\pi-x}\pi\int_0^x y\,\dd y
  +\frac x\pi\int_x^\pi(\pi-y)\,\dd y\\
 &=\frac{x(\pi-x)}2.
\end{align*}
常源的谱系数为
\[
 b_n=\langle1,e_n\rangle
 =\sqrt{2/\pi}\,\frac{1-(-1)^n}{n},
 \qquad \frac{b_n}{n^2}=\frac{c_n}{2}.
\]
所以 $A_D^{-1}1=\sum b_nn^{-2}e_n=f/2$，三个解法完全相同。
Neumann 则不同：若 $-u''=1$ 且两端导数为零，积分给 $0=-u'(\pi)+u'(0)=\pi$，产生矛盾。
这是源没有与核 $\operatorname{span}\{1\}$ 正交的 Fredholm 相容性失败。

记 $K=A_D^{-1}$，它的 SVD 在本题中就是
\[
 Ku=\sum_{n\ge1}s_n\langle u,e_n\rangle e_n,
 \qquad s_n=n^{-2}.
\]
左右奇异向量相同，是因为 $K$ 正、自伴。
令 $K_N=\sum_{n=1}^N n^{-2}\langle\cdot,e_n\rangle e_n$，则
\[
 \|K-K_N\|=\frac1{(N+1)^2}\longrightarrow0.
\]
所以 $K$ 是有限秩算子的算子范数极限，因而紧。
而任何秩不超过 $N$ 的 $F$，在 $\operatorname{span}\{e_1,\ldots,e_{N+1}\}$ 上都有非零核向量 $w$，取 $\|w\|=1$ 得
\[
 \|K-F\|\ge\|Kw\|\ge\frac1{(N+1)^2}.
\]
所以 $K_N$ 同时给最佳算子范数有限秩逼近。
进一步，
\[
 \|K\|=1,\qquad \|K\|_{\mathrm{HS}}^2=\sum n^{-4}=\frac{\pi^4}{90},
 \qquad \operatorname{tr}K=\sum n^{-2}=\frac{\pi^2}{6}.
\]
最后一个数也等于 $\int_0^\pi G_D(0;x,x)\,\dd x$。
\end{exbox}

\begin{exbox}[同一初态的幺正演化、热演化与源驱动]
第八步得到恒源的静态响应 $f/2$。现在回到同一个初态，用已经确定的 $A_D$ 分别求
\[
 \ii\partial_t\psi=A_D\psi,\qquad
 \partial_t w=-A_Dw,
 \qquad \psi(0)=w(0)=f.
\]

\textbf{解。} 谱演算给出完整解
\begin{align*}
 \psi(t,x)&=\frac8\pi\sum_{n\text{ 奇}}\frac{e^{-\ii tn^2}}{n^3}\sin(nx),\\
 w(t,x)&=\frac8\pi\sum_{n\text{ 奇}}\frac{e^{-tn^2}}{n^3}\sin(nx),\qquad t\ge0.
\end{align*}
第一族 $U_D(t)=e^{-\ii tA_D}$ 是强连续幺正群：相位相乘给群律，$|e^{-\ii tn^2}|=1$ 给范数保持，支配收敛给强连续性。
在 $\Dom(A_D)$ 上逐系数求导并以 $\sum n^4|c_n|^2<\infty$ 控制尾项，得到 $\ii\psi'=A_D\psi$。
本题初态正好满足这一要求，且 $\psi'(0)=-2\ii$ 在 $L^2$ 意义成立。
但不能再要求二阶强时间导数，因为 $f\notin\Dom(A_D^2)$。

热半群的每个模态衰减，满足
\[
 \|w(t)\|^2=\frac{32}\pi\sum_{n\text{ 奇}}\frac{e^{-2tn^2}}{n^6},
 \qquad \|w(t)\|\le e^{-t}\|f\|.
\]
对 $t>0$ 和每个整数 $m\ge0$，$\sum n^{4m}e^{-2tn^2}|c_n|^2<\infty$，故热解进入全部 $\Dom(A_D^m)$。
Schr\"odinger 演化只改相位，不产生这种平滑作用。

若再加入同一个恒定源 $1$，Duhamel 公式要求积分每个模态的传播因子。
对 Schr\"odinger 方程，积分项为 $-\ii b_n\int_0^t e^{-\ii(t-s)n^2}\,\dd s$；
对热方程则为 $b_n\int_0^t e^{-(t-s)n^2}\,\dd s$。分别完成这两个积分，得到
\begin{align*}
 \ii\partial_t\psi_s&=A_D\psi_s+1,&
 \psi_s(t)&=\sum_n\left[e^{-\ii tn^2}c_n
 -\frac{1-e^{-\ii tn^2}}{n^2}b_n\right]e_n,\\
 \partial_t w_s&=-A_Dw_s+1,&
 w_s(t)&=\sum_n\left[e^{-tn^2}c_n
 +\frac{1-e^{-tn^2}}{n^2}b_n\right]e_n.
\end{align*}
由于 $b_n/n^2=c_n/2$，热源解进一步化成
\[
 w_s(t)=\frac f2+\frac12e^{-tA_D}f\longrightarrow\frac f2.
\]
极限正是第八步算出的静态响应。
带源 Schr\"odinger 解的静态特解则为 $-f/2$，其余部分作幺正振荡；符号不同来自两个方程中源的位置。

对 Neumann 热演化，常数谱投影不衰减：
\[
 e^{-tA_N}f=d_0\eta_0+\sum_{n\ge1}e^{-tn^2}d_n\eta_n
 \longrightarrow\frac{\pi^2}{6}.
\]
再加常源 $1$ 后，零模随时间增长，解为 $e^{-tA_N}f+t$，不会趋于稳态。
零模既决定静态相容性，也决定长期演化。
\end{exbox}

\begin{exbox}[有限模计算的最佳性与误差]
第九步的完整演化包含无穷多个模态。实际计算令 $V_N=\operatorname{span}\{e_1,\ldots,e_N\}$，只保留前 $N$ 个模态。
需要分别回答：这是否是最佳近似，误差在哪种范数中小，时间演化会不会放大误差？

\textbf{解。} 正交投影为
\[
 P_Nu=\sum_{n=1}^N\langle u,e_n\rangle e_n.
\]
对任意 $v\in V_N$，正交分解给
\[
 \|u-v\|^2=\|u-P_Nu\|^2+\|P_Nu-v\|^2.
\]
所以 $P_Nu$ 是唯一最佳 $L^2$ 近似，而且 $P_N^2=P_N=P_N^*$。
本题三个不同强度的精确尾误差为
\begin{align*}
 \|f-P_Nf\|^2&=\frac{32}\pi\sum_{\substack{n>N\\n\text{ 奇}}}n^{-6}
 \le\frac{32}{5\pi N^5},\\
 q_D(f-P_Nf,f-P_Nf)&=\frac{32}\pi\sum_{\substack{n>N\\n\text{ 奇}}}n^{-4}
 \le\frac{32}{3\pi N^3},\\
 \|A_D(f-P_Nf)\|^2&=\frac{32}\pi\sum_{\substack{n>N\\n\text{ 奇}}}n^{-2}
 \le\frac{32}{\pi N}.
\end{align*}
上界使用递减函数的积分估计 $\sum_{n>N}n^{-p}\le\int_N^\infty x^{-p}\,\dd x$，$N\ge1$。
同一次截断在 $L^2$、能量范数和图范数中有不同收敛速度；不能用第一行代替最后一行。

例如只保留 $n=1$ 时
\[
 P_1f=\frac8\pi\sin x,\qquad
 \|f-P_1f\|^2=\frac{\pi^5}{30}-\frac{32}\pi,
 \qquad
 \frac{\|f-P_1f\|^2}{\|f\|^2}=1-\frac{960}{\pi^6}.
\]
一般地，要求相对 $L^2$ 误差不超过 $\varepsilon$，一个可直接采用的充分条件是
\[
 N\ge\left(\frac{192}{\pi^6\varepsilon^2}\right)^{1/5}.
\]
这是保守上界；由于偶数模态恰为零，实际误差由奇数尾和给出。

谱投影与 $A_D$、$U_D(t)$、$e^{-tA_D}$ 对易，因此有限模解就是 $\psi_N(t)=U_D(t)P_Nf$ 与 $w_N(t)=e^{-tA_D}P_Nf$。
误差满足
\[
 \|\psi(t)-\psi_N(t)\|=\|f-P_Nf\|,
 \qquad
 \|w(t)-w_N(t)\|\le e^{-t(N+1)^2}\|f-P_Nf\|.
\]
幺正演化保持尾误差，热演化更快压低高频尾模态。
对带恒源的热解，利用 $w_s=f/2+e^{-tA_D}f/2$，得到
\[
 \|w_s(t)-P_Nw_s(t)\|
 \le\frac{1+e^{-t(N+1)^2}}2\|f-P_Nf\|.
\]
静态解 $u=f/2$ 的误差恰为初态误差的一半。

这些误差估计使用的仍是初态系数 $c_n$ 和源系数 $b_n$。
边界条件先决定本征系，初态的算子域决定可求的时间导数，所选范数再决定截断收敛的速度；三者共同决定有限模解能近似什么。
\end{exbox}
